% ----------------------------------------------------------------------------
%  19. Conclusion
% ----------------------------------------------------------------------------
\section{Conclusion}
\label{sec:conclusion}

\subsection{What was implemented}

\begin{itemize}
  \item A scheduler with two pluggable axes: ordering and
        matching.
  \item Three matchers (Greedy, FairGreedy, Hungarian).
  \item Seven priority policies subsuming the legacy formula.
  \item A multi-objective scoring step.
  \item A six-mode benchmark suite over a 100-patient,
        10-worker ward.
\end{itemize}

\subsection{What the experiments showed}

\begin{itemize}
  \item Ordering moves more than matching.
  \item ``Optimal per tick'' (Hungarian) is the wrong invariant
        when many critical tasks share the same priority --- it
        regresses worst-case CRIT timing by $4\times$ vs.\ Greedy.
  \item No algorithm change creates capacity. The ward is
        $\sim 5\times$ over-loaded, so 20{,}210 minutes of overtime
        appear under every strategy.
\end{itemize}

\subsection{Closing note}

\begin{itemize}
  \item The contribution is not a new algorithm; it is the
        disciplined separation of ``what task next?'' from
        ``who does it?'', and a benchmark surface that lets a
        single observed difference be attributed to a single
        choice.
  \item The headline finding is humbling --- the simplest matcher
        won --- but that's exactly the kind of finding the suite
        was built to surface.
\end{itemize}
